\chapter{Memory Subsystem}

\section{Overview}

The Memory Subsystem manages all load and store operations generated by the processor.
Rather than accessing memory directly from the execution stage, Sargantana uses dedicated
hardware structures to support efficient memory scheduling, dependency tracking and store
buffering.

The primary RTL modules involved in memory operations are:

\begin{rtlrefs}
  \rtlitem{rtl/datapath/rtl/exe_stage/rtl/mem_unit.sv}
  \rtlitem{rtl/datapath/rtl/exe_stage/rtl/load_store_queue.sv}
  \rtlitem{rtl/datapath/rtl/exe_stage/rtl/store_buffer.sv}
  \rtlitem{rtl/datapath/rtl/exe_stage/rtl/vagu.sv}
  \rtlitem{rtl/datapath/rtl/interface_csr/rtl/csr_interface.sv}
\end{rtlrefs}

These modules collectively implement the processor's memory execution pipeline.

\section{Load Operations}

When a load instruction reaches the Execution stage, the effective address is first computed by
the Address Generation Unit (\rtl{vagu.sv}). The request is then inserted into the Load/Store
Queue, which manages pending memory operations.

The Memory Unit (\rtl{mem\_unit.sv}) communicates with the cache hierarchy and retrieves the
requested data. Once the data becomes available, it is forwarded toward the Writeback stage,
where it updates the corresponding physical register.

This separation allows multiple outstanding memory requests to be tracked without stalling the
entire pipeline.

\section{Store Operations}

Store instructions follow a similar flow but differ in that they write data to memory rather than
returning a result to the register file.

Before a store is committed to memory, it is temporarily held in the Store Buffer:

\begin{center}
\rtl{rtl/datapath/rtl/exe\_stage/rtl/store\_buffer.sv}
\end{center}

The Store Buffer allows execution to continue while pending store operations are safely queued.
Once it is safe to update memory, buffered store entries are committed in program order.

\section{Address Generation}

Effective memory addresses are generated by the dedicated Address Generation Unit.

RTL implementation:

\begin{center}
\rtl{rtl/datapath/rtl/exe\_stage/rtl/vagu.sv}
\end{center}

The Address Generation Unit combines the base register value with the instruction offset to
compute the final virtual address used by both load and store operations.

Separating address generation into its own module improves modularity and enables future
support for more advanced addressing modes.

\section{Cache Interface}

The Memory Unit serves as the interface between the execution pipeline and the processor's
cache hierarchy. It issues memory requests generated by the Load/Store Queue and receives
the corresponding responses.

The cache interface logic is primarily coordinated by:

\begin{center}
\rtl{rtl/datapath/rtl/exe\_stage/rtl/mem\_unit.sv}
\end{center}

and integrated into the top-level datapath through:

\begin{center}
\rtl{rtl/datapath/rtl/datapath.sv}
\end{center}

This organization isolates cache communication from the rest of the execution logic while
allowing memory latency to be managed independently.

\section{Memory Ordering}

Correct program execution requires that memory operations observe architectural ordering
rules. Sargantana enforces this through the interaction of the Load/Store Queue, Store Buffer
and scoreboard structures.

The relevant RTL modules include:

\begin{rtlrefs}
  \rtlitem{load_store_queue.sv}
  \rtlitem{store_buffer.sv}
  \rtlitem{pending_mem_req_queue.sv}
  \rtlitem{score_board_scalar.sv}
\end{rtlrefs}

Together, these modules track outstanding memory requests, preserve dependencies between
loads and stores, and ensure that memory operations are completed in a correct and consistent
order before instructions are retired.
